%% Question bank chapter file - generated by bank_tools.py build (do not edit while the bank is still being built:
%% it is regenerated from the per-test data in Question Bank\data). Every question is one bankq block:
%%   \begin{bankq}[id=..., source={...}, paper=.., qno=.., type=..., subtopic={...}, answer={...}, answer_src=key|solved,
%%                 working={...}, difficulty=1-10, diff_parts={K S I C F}, diff_note={...},
%%                 group=<passage id or empty>, also_in={duplicate copies in other tests}]
%% difficulty = min(10, K+S+I+C+F-1): K concepts 1-3, S steps 1-4, I insight 0-2, C computation 0-2, F format 0-1
%%   ...question LaTeX, options, figures (\bfig / \bfigi -> img/)...
%%   \end{bankq}
%% Passages shared by several questions are bankpassage blocks; their questions carry group=<passage id>.
\documentclass[10pt]{article}
\usepackage[a4paper,margin=1.8cm]{geometry}
\usepackage[T1]{fontenc}
\usepackage{amsmath,amssymb,mathtools}
\usepackage{graphicx}
\usepackage[export]{adjustbox}
\usepackage{xcolor,needspace,keyval}
\usepackage{enumitem}
\usepackage{booktabs,array,tabularx}
\usepackage[hidelinks]{hyperref}
\setlength{\parindent}{0pt}
\setlength{\parskip}{3pt}
\makeatletter
\define@key{bq}{id}{\def\bq@id{#1}}\define@key{bq}{source}{\def\bq@src{#1}}\define@key{bq}{paper}{\def\bq@paper{#1}}
\define@key{bq}{qno}{\def\bq@qno{#1}}\define@key{bq}{type}{\def\bq@type{#1}}\define@key{bq}{subtopic}{\def\bq@sub{#1}}
\define@key{bq}{answer}{\def\bq@ans{#1}}\define@key{bq}{answer_src}{\def\bq@asrc{#1}}\define@key{bq}{working}{\def\bq@work{#1}}
\define@key{bq}{group}{\def\bq@group{#1}}\define@key{bq}{also_in}{\def\bq@also{#1}}
\define@key{bq}{difficulty}{\def\bq@diff{#1}}\define@key{bq}{diff_parts}{\def\bq@dparts{#1}}\define@key{bq}{diff_note}{\def\bq@dnote{#1}}
\newenvironment{bankq}[1][]{\def\bq@work{}\def\bq@also{}\def\bq@group{}\def\bq@diff{?}\def\bq@dparts{}\def\bq@dnote{}\setkeys{bq}{#1}%
  \par\needspace{4\baselineskip}\medskip\noindent{\small\textbf{\bq@id}\hfill\color{gray}[\bq@src, Paper \bq@paper, Q\bq@qno]
  \bq@type\ $\cdot$ \bq@sub\ $\cdot$ \textbf{D\bq@diff}}\par\nopagebreak\relax}
  {\par{\small\color{teal}Answer: \textbf{\bq@ans}\ifx\bq@work\@empty\else\ (solved: \bq@work)\fi
  \ifx\bq@also\@empty\else\par\color{gray}Also in: \bq@also\fi}\par}
\define@key{bp}{id}{\def\bp@id{#1}}\define@key{bp}{source}{\def\bp@src{#1}}\define@key{bp}{questions}{\def\bp@qs{#1}}
\newenvironment{bankpassage}[1][]{\setkeys{bp}{#1}\par\needspace{4\baselineskip}\bigskip\noindent
  {\small\textbf{Passage \bp@id}\hfill\color{gray}[\bp@src; questions \bp@qs]}\par\nopagebreak\relax\itshape}{\par}
\makeatother
\newcommand{\bfig}[2]{\par\smallskip{\centering\includegraphics[width=#1pt,max width=0.62\linewidth,max height=0.32\textheight]{img/#2}\par}\smallskip}
\newcommand{\bfigi}[2]{\includegraphics[width=#1pt,max width=\linewidth,max height=0.32\textheight,valign=c]{img/#2}}
\newcommand{\bmissing}[1]{\textit{[figure missing in the source conversion: \detokenize{#1}]}}
\newenvironment{qoptions}{\begin{itemize}[leftmargin=1.8em,itemsep=1pt,topsep=2pt,label={}]}{\end{itemize}}
\newcommand{\opt}[2]{\item[(#1)] #2}
\providecommand{\fwstart}{}\providecommand{\fwend}{}\providecommand{\qfullsplit}{}
\newcolumntype{L}[1]{>{\raggedright\arraybackslash}p{#1}}
\newcolumntype{C}[1]{>{\centering\arraybackslash}p{#1}}
\newcolumntype{R}[1]{>{\raggedleft\arraybackslash}p{#1}}
\providecommand{\sectionhdr}[1]{}\providecommand{\subsectionhdr}[1]{}


\begin{document}
{\Large\bfseries Straight Line \& Pair --- Question Bank}\par\smallskip
{\small Mathematics chapter M15. Questions from all converted test papers, grouped by source series. Answers: printed key where it matches the paper, otherwise solved (answer\_src=solved, with working). Exact repeats across papers appear once, with also\_in listing the other copies.}\par\medskip\hrule

\section*{Score-I papers}
\begin{bankpassage}[id=SI12-P2-MATH-S1-G5-6, source={Score-I - Test 12 (09 Mar 2022)}, questions={5-6}]
Four mirrors are placed along the lines $y=0$, $x=1$, $y=6$ and $x=-10$, a ray of line is emerging from origin, it reflects from the mirrors and exits as soon as it reaches the intersection of any two mirrors.\\
Column 1 : Denotes slope of incident ray\\
Column 2 : Denotes number of reflections before it exits\\
Column 3 : Denotes length of path covered before it exists
\begin{center}\footnotesize
\begin{tabular}{|p{0.6cm}|p{1.4cm}|p{0.6cm}|p{1.4cm}|p{0.5cm}|p{1.8cm}|}
\hline
 & \centering Column 1 & & \centering Column 2 & & \centering Column 3\tabularnewline
\hline
(I) & $m=1$ & (i) & 8 & (P) & $10\sqrt{10}$\\
(II) & $m=2$ & (ii) & 6 & (Q) & $12\sqrt5$\\
(III) & $m=3$ & (iii) & 2 & (R) & $12\sqrt2$\\
(IV) & $m=-3$ & (iv) & 4 & (S) & $12\sqrt{10}$\\
\hline
\end{tabular}
\end{center}
\end{bankpassage}
\begin{bankq}[id=SI12-P2-MATH-S1-Q05, source={Score-I - Test 12 (09 Mar 2022)}, paper=2, qno=5, type=single, subtopic={Forms of lines}, answer={B}, answer_src=key, difficulty=2, diff_parts={K1 S2 I0 C0 F0}, diff_note={concepts: forms of straight lines | steps: line for item I; line for item II; line for item III; line for item IV and match}, group=SI12-P2-MATH-S1-G5-6]
Identify the correct combination ?
\begin{qoptions}
\opt{A}{(I) (iii) (S)}
\opt{B}{(II) (iv) (Q)}
\opt{C}{(III) (i) (P)}
\opt{D}{(IV) (iv) (Q)}
\end{qoptions}
\end{bankq}
\begin{bankq}[id=SI12-P2-MATH-S1-Q06, source={Score-I - Test 12 (09 Mar 2022)}, paper=2, qno=6, type=single, subtopic={Forms of lines}, answer={D}, answer_src=key, difficulty=2, diff_parts={K1 S2 I0 C0 F0}, diff_note={concepts: forms of straight lines | steps: line for item I; line for item II; line for item III; line for item IV and match}, group=SI12-P2-MATH-S1-G5-6]
Identify the correct combination ?
\begin{qoptions}
\opt{A}{(I) (iii) (Q)}
\opt{B}{(II) (iv) (P)}
\opt{C}{(III) (ii) (R)}
\opt{D}{(IV) (iv) (P)}
\end{qoptions}
\end{bankq}

\section*{TOAS papers}
\begin{bankq}[id=TOAS01-P2-MATH-S1-Q02, source={TOAS - Test 01 (30 Jan 2022)}, paper=2, qno=2, type=single, subtopic={Angle between lines}, answer={C}, answer_src=key, difficulty=4, diff_parts={K2 S2 I1 C0 F0}, diff_note={concepts: angle conditions, slope | steps: angle at P between B and the sides; locate P; slope of BP | I: BP bisects the angle APC}]
A point $P$ is taken outside $\Delta ABC$ where $(B(1,\sqrt3)$, $A(0,0)$ and $C(2,0))$, but inside the acute angle $BAC$, such that $\angle APC=\dfrac\pi6$ and $\angle BPA=\dfrac\pi{12}$. Slope of the line $BP$ is
\begin{qoptions}
\opt{A}{$\sqrt3$}
\opt{B}{$-\sqrt3$}
\opt{C}{$\dfrac1{\sqrt3}$}
\opt{D}{$\dfrac{-1}{\sqrt3}$}
\end{qoptions}
\end{bankq}
\begin{bankq}[id=TOAS01-P2-MATH-S1-Q05, source={TOAS - Test 01 (30 Jan 2022)}, paper=2, qno=5, type=single, subtopic={Locus}, answer={A}, answer_src=key, difficulty=2, diff_parts={K1 S1 I1 C0 F0}, diff_note={concepts: perimeter of the offset curve | steps: parallel offset adds the circular arcs 2 pi 3 | I: offset path}]
A man moves outside a park $(ABCDE)$ in such a way that he always remains at a distance of 3m from the wall (as shown in figure), then total distance travelled for a complete round is -
\bfig{77.9}{TOAS01_p2_math_s1_q05_1.png}
\begin{qoptions}
\opt{A}{$(200+6\pi)$ m}
\opt{B}{$(200+8\pi)$ m}
\opt{C}{$(200+10\pi)$ m}
\opt{D}{$(200+4\pi)$ m}
\end{qoptions}
\end{bankq}
\begin{bankq}[id=TOAS06-P1-MATH-S1-Q04, source={TOAS - Test 06 (17 Apr 2022)}, paper=1, qno=4, type=single, subtopic={Forms of lines}, answer={B}, answer_src=key, difficulty=4, diff_parts={K2 S2 I1 C0 F0}, diff_note={concepts: angle bisector, perpendicular bisector | steps: reflect A in the bisector to lie on BC; line BC through two points; slope m | I: reflection in the bisector lies on BC}]
In triangle ABC, vertex A is $(1,1)$ and internal angle bisector $x-y=2$ of angle B meets the perpendicular bisector of side AC at $(4,2)$. If equation of side AC is $2x+y=3$ and slope of side BC is $m$, then $|4m|$ is equal to
\begin{qoptions}
\opt{A}{4}
\opt{B}{3}
\opt{C}{6}
\opt{D}{7}
\end{qoptions}
\end{bankq}
\begin{bankq}[id=TOAS06-P1-MATH-S2-Q07, source={TOAS - Test 06 (17 Apr 2022)}, paper=1, qno=7, type=multi, subtopic={Pair of lines}, answer={C,D}, answer_src=key, difficulty=5, diff_parts={K2 S2 I1 C0 F1}, diff_note={concepts: pair of lines, quadrilateral | steps: angle bisectors of the pairs; shape of the quadrilateral; area as a function of K; range | I: the second pair has the same bisectors}]
Pairs of lines $\sqrt3(x-2)^2-4(x-2)(y-2)+\sqrt3(y-2)^2=0$ \quad and \quad $\sqrt3x^2-Kxy+\sqrt3y^2=0$; $K\in[2\sqrt3,\,4]$ intersect to form a quadrilateral of area $A_K$.
\begin{qoptions}
\opt{A}{Range of $A_K\in[1,\ \sqrt3]$}
\opt{B}{Quadrilateral is rhombus $\forall\ K\in[2\sqrt3,\ 4]$}
\opt{C}{Range of $A_K\in[0,\ 8-4\sqrt3]$}
\opt{D}{Quadrilateral is kite $\forall\ K\in[2\sqrt3,\ 4]$}
\end{qoptions}
\end{bankq}
\begin{bankq}[id=TOAS06-P1-MATH-S2-Q09, source={TOAS - Test 06 (17 Apr 2022)}, paper=1, qno=9, type=multi, subtopic={Forms of lines}, answer={A,C}, answer_src=key, difficulty=5, diff_parts={K2 S2 I1 C0 F1}, diff_note={concepts: incentre, altitude | steps: incentre on the altitude from A fixes alpha; line AC by reflection; vertices and area | I: reflection of AB in the bisector}]
Let two sides AB and BC of a $\Delta$ABC be represented by $x-7y+3=0$ and $x+y+3=0$ respectively. Incentre I $(-3-3\alpha,\ \alpha)$, $\alpha\in R$ of the triangle lies on the altitude $x-y-1=0$ through A then
\begin{qoptions}
\opt{A}{The equation of a line representing side AC is $7x-y=11$}
\opt{B}{The equation of a line representing side AC is $7x+y=1$}
\opt{C}{The area of $\Delta$ABC is $\dfrac{32}{3}$ sq. unit}
\opt{D}{The area of $\Delta$ABC is $\dfrac{32}{5}$ sq.unit}
\end{qoptions}
\end{bankq}
\begin{bankpassage}[id=TOAS07-P2-MATH-S2-G7-8, source={TOAS - Test 07 (01 May 2022)}, questions={7-8}]
A line passes through the point $P(2,3)$ makes an angle $\theta$ with positive direction of $x$-axis and it meets the lines represented by $x^2-2xy-y^2=0$ at the points $A$ and $B$.\\
On the basis of above information, answer the following questions :
\end{bankpassage}
\begin{bankq}[id=TOAS07-P2-MATH-S2-Q07, source={TOAS - Test 07 (01 May 2022)}, paper=2, qno=7, type=single, subtopic={Pair of lines}, answer={D}, answer_src=key, difficulty=3, diff_parts={K1 S2 I0 C1 F0}, diff_note={concepts: derivative of a sum of lengths | steps: express PA and PB in terms of theta; differentiate; evaluate at pi/4}, group=TOAS07-P2-MATH-S2-G7-8]
$\left(\dfrac{d}{d\theta}(PA+PB)\right)_{\theta=\frac\pi4}$ is equal to -
\begin{qoptions}
\opt{A}{0}
\opt{B}{$\sqrt2$}
\opt{C}{$4\sqrt2$}
\opt{D}{$8\sqrt2$}
\end{qoptions}
\end{bankq}
\begin{bankq}[id=TOAS07-P2-MATH-S2-Q08, source={TOAS - Test 07 (01 May 2022)}, paper=2, qno=8, type=single, subtopic={Pair of lines}, answer={C}, answer_src=key, difficulty=1, diff_parts={K1 S1 I0 C0 F0}, diff_note={concepts: product of lengths | steps: PA times PB as a function of theta; solve for 17}, group=TOAS07-P2-MATH-S2-G7-8]
If $PA.PB=17$, then $\theta$ can be equal to -
\begin{qoptions}
\opt{A}{$\dfrac\pi3$}
\opt{B}{$\dfrac\pi6$}
\opt{C}{$\dfrac\pi2$}
\opt{D}{$\dfrac{2\pi}{3}$}
\end{qoptions}
\end{bankq}
\begin{bankpassage}[id=TOAS09-P1-MATH-S3-G1-2, source={TOAS - Test 09 (30 May 2022)}, questions={1-2}]
A Rod of length 10cm slides between two lines, x-axis \& $y=-\sqrt3\,x$. If locus of mid point of rod is curve `S'.
\end{bankpassage}
\begin{bankq}[id=TOAS09-P1-MATH-S3-Q01, source={TOAS - Test 09 (30 May 2022)}, paper=1, qno=1, type=numerical, subtopic={Locus}, answer={10.00}, answer_src=key, difficulty=2, diff_parts={K1 S1 I0 C0 F1}, diff_note={concepts: locus from the passage | steps: derive the equation of S; x-intercept}, group=TOAS09-P1-MATH-S3-G1-2]
x-intercept of S is
\end{bankq}
\begin{bankq}[id=TOAS09-P1-MATH-S3-Q02, source={TOAS - Test 09 (30 May 2022)}, paper=1, qno=2, type=numerical, subtopic={Locus}, answer={15.00}, answer_src=key, difficulty=3, diff_parts={K1 S2 I0 C0 F1}, diff_note={concepts: locus from the passage | steps: parametrise S; maximise distance from the origin; k sqrt3}, group=TOAS09-P1-MATH-S3-G1-2]
If maximum distance between origin and curve S is equal to $k$ then value of $k\sqrt3$ is
\end{bankq}
\begin{bankq}[id=TOAS09-P2-MATH-S1-Q02, source={TOAS - Test 09 (30 May 2022)}, paper=2, qno=2, type=single, subtopic={Forms of lines}, answer={B}, answer_src=key, difficulty=4, diff_parts={K1 S2 I1 C1 F0}, diff_note={concepts: point inside a triangle formed by lines | steps: vertices of the triangle; inequality from each line; combine the conditions on a | I: sketch the region for the point on a curve}]
The complete set of values the parameter `a' so that the point $P\left(a,\ \dfrac{1}{1+a^2}\right)$ doesn't lie outside the triangle formed by the lines $15y=x+1$, $78y=118-23x$ and $y+2=0$ is -
\begin{qoptions}
\opt{A}{$(0,5)$}
\opt{B}{$[2,5]$}
\opt{C}{$(1,5)$}
\opt{D}{$[0,2]$}
\end{qoptions}
\end{bankq}
\begin{bankq}[id=TOAS09-P2-MATH-S1-Q12, source={TOAS - Test 09 (30 May 2022)}, paper=2, qno=12, type=single, subtopic={Pair of lines}, answer={D}, answer_src=key, difficulty=4, diff_parts={K2 S2 I0 C1 F0}, diff_note={concepts: pair of lines, parallelogram area | steps: factor the pair of lines; vertices from the diagonal; area of the parallelogram}]
If two adjacent sides of a parallelogram are given by $2x^2+2y^2-5xy=0$ and $5x+2y-1=0$ is one of its diagonal, then area of parallelogram is $\dfrac1k$ sq. units, where $k$ is
\begin{qoptions}
\opt{A}{16}
\opt{B}{72}
\opt{C}{18}
\opt{D}{36}
\end{qoptions}
\end{bankq}
\end{document}
